\documentclass[a4paper,10pt]{article}
\usepackage{lamia-pedagogy}

\begin{document}

% =========================================================================
% PAGE 1 : Texte Support & Démarche de Lecture
% =========================================================================
\ficheHeader{1 Heure}{\makebox[2.5cm]{\dotfill}}{1\textsuperscript{ère} AM}{2026 / 2027}

\vspace{-3pt}
\competenceBanner{\faBookReader\ Projet 01 / Séquence 01 — Séance 03 : Compréhension de l'Écrit (01)}

\vspace{-2pt}
\modaliteSupportBar{Lecture silencieuse puis collective}{Texte explicatif \guillemotleft\ Des gestes simples pour une bonne hygiène corporelle \guillemotright}

\vspace{1pt}

% Texte Support
\begin{tcolorbox}[
  enhanced,
  colback=primarylight!40,
  colframe=primary,
  arc=2mm,boxrule=0.8pt,
  top=5pt,bottom=5pt,left=8pt,right=8pt,
  title={\small\textbf{\textcolor{white}{\faFile*\ Texte : Des gestes simples pour une bonne hygiène corporelle}}},
  coltitle=white,colbacktitle=primary
]
{\small
À la maison et à l'école, les adultes demandent souvent aux enfants de se laver les mains. Mais pourquoi ce geste est-il si important ?

\textbf{Le lavage des mains} consiste à \underline{se débarrasser} des saletés et des germes invisibles. Ces microbes passent facilement d'un objet contaminé (poignée de porte, jouet, sol) aux mains d'une personne. Si on ne se lave pas les mains avant de manger ou après avoir utilisé les toilettes, les microbes pénètrent dans notre corps et provoquent des maladies parfois graves.

Pour bien se laver les mains, il ne suffit pas de les mouiller : il faut obligatoirement utiliser de l'eau et du savon, bien frotter pendant au moins 30 secondes, puis les sécher avec une serviette propre. Avoir des mains propres est donc un geste simple qui nous permet d'éviter d'être malades et de protéger la santé des autres.

\hfill\textit{\footnotesize D'après l'Organisation Mondiale de la Santé (OMS)}
}
\end{tcolorbox}

\vspace{3pt}

% Structure du texte
\begin{tcolorbox}[
  enhanced,
  colback=white,
  colframe=cardborder,
  arc=2mm,boxrule=0.7pt,
  top=4pt,bottom=4pt,left=7pt,right=7pt,
  title={\small\textbf{\textcolor{primary}{\faLayerGroup\ Grille de lecture \& Hypothèses de sens}}},
  coltitle=primary,colbacktitle=primarylight,titlerule=0.5pt
]
{\footnotesize
\begin{tabularx}{\linewidth}{|X|X|X|}
\hline
\rowcolor{primarylight}\textbf{1. Phase de questionnement} & \textbf{2. Phase d'explication} & \textbf{3. Phase de conclusion} \\
\hline
Introduction du thème et question de départ (\textit{Pourquoi ce geste est-il important ?}). &
Définition du lavage des mains, rôle des microbes et méthode de nettoyage. &
Bilan sur les bienfaits de l'hygiène pour la santé individuelle et collective. \\
\hline
\end{tabularx}
}
\end{tcolorbox}

\vspace{3pt}

% Objectifs d'apprentissage
\begin{tcolorbox}[
  enhanced,
  colback=secondarylight!50,
  colframe=secondary,
  arc=1.5mm,leftrule=3pt,rightrule=0.5pt,toprule=0.5pt,bottomrule=0.5pt,
  top=3pt,bottom=3pt,left=6pt,right=6pt
]
{\footnotesize
\textbf{\faBullseye\ Objectifs visés :} Identifier les caractéristiques d'un texte explicatif, repérer une définition et justifier une réponse à partir du texte.
}
\end{tcolorbox}

\newpage

% =========================================================================
% PAGE 2 : Fiche d'Exploitation et d'Évaluation de l'Élève
% =========================================================================
\ficheEleveHeader

\vspace{-3pt}
\competenceBanner{\faPenFancy\ Questionnaire de Compréhension de l'Écrit}

\vspace{2pt}

% ACTIVITÉ 1 : Choix multiple
\begin{tcolorbox}[
  enhanced,
  colback=white,
  colframe=cardborder,
  arc=2mm,boxrule=0.8pt,
  top=4pt,bottom=4pt,left=8pt,right=8pt,
  title={\small\textbf{\textcolor{primary}{\faCheckSquare\ Activité 01 : Choisis la bonne réponse}} \hfill \badge[primary]{Repérage}},
  coltitle=primary,colbacktitle=primarylight,titlerule=0.5pt,titlerule style=cardborder
]
{\small
\begin{enumerate}[label=\textbf{\arabic*.},leftmargin=*,noitemsep,topsep=2pt]
  \item Dans ce texte, l'auteur parle :
    \cardOptThree{des maladies}{du sport}{du lavage des mains}
  \item Pour bien se laver les mains, on utilise :
    \cardOptThree{la brosse à dents}{l'eau et le savon}{le shampoing}
  \item Le mot souligné \guillemotleft\ \textbf{se débarrasser} \guillemotright\ veut dire :
    \cardOptThree{laver}{protéger}{éliminer}
\end{enumerate}
}
\end{tcolorbox}

\vspace{2pt}

% ACTIVITÉ 2 : Vrai ou Faux
\begin{tcolorbox}[
  enhanced,
  colback=white,
  colframe=cardborder,
  arc=2mm,boxrule=0.8pt,
  top=4pt,bottom=4pt,left=8pt,right=8pt,
  title={\small\textbf{\textcolor{primary}{\faCheckDouble\ Activité 02 : Réponds par Vrai ou Faux}} \hfill \badge[secondary]{Validation}},
  coltitle=primary,colbacktitle=primarylight,titlerule=0.5pt,titlerule style=cardborder
]
{\small
\begin{enumerate}[label=\textbf{\alph*)},leftmargin=*,noitemsep,topsep=2pt]
  \item Les mains propres causent des maladies graves. \dotfill \cardOptTwo{Vrai}{Faux}
  \item Le lavage des mains consiste à éliminer les saletés et les microbes. \dotfill \cardOptTwo{Vrai}{Faux}
  \item Les microbes passent facilement d'un objet contaminé aux mains. \dotfill \cardOptTwo{Vrai}{Faux}
  \item Il suffit de mouiller les mains avec de l'eau seule pour être propre. \dotfill \cardOptTwo{Vrai}{Faux}
\end{enumerate}
}
\end{tcolorbox}

\vspace{2pt}

% ACTIVITÉ 3 : Compréhension & Définition
\begin{tcolorbox}[
  enhanced,
  colback=white,
  colframe=cardborder,
  arc=2mm,boxrule=0.8pt,
  top=4pt,bottom=4pt,left=8pt,right=8pt,
  title={\small\textbf{\textcolor{primary}{\faQuestionCircle\ Activité 03 : Réponds aux questions suivantes}} \hfill \badge[accent]{Analyse}},
  coltitle=primary,colbacktitle=primarylight,titlerule=0.5pt,titlerule style=cardborder
]
{\small
\begin{enumerate}[label=\textbf{\arabic*.},leftmargin=*,noitemsep,topsep=2pt]
  \item Relève dans le texte la phrase qui définit le lavage des mains :\\
    \drule{\linewidth}\\[4pt]
    \drule{\linewidth}
  \item Que risquent les personnes qui n'ont pas une bonne hygiène des mains ?\\
    \drule{\linewidth}
\end{enumerate}
}
\end{tcolorbox}

\vspace{2pt}

% ACTIVITÉ 4 : Synthèse
\begin{tcolorbox}[
  enhanced,
  colback=white,
  colframe=cardborder,
  arc=2mm,boxrule=0.8pt,
  top=4pt,bottom=4pt,left=8pt,right=8pt,
  title={\small\textbf{\textcolor{primary}{\faEdit\ Activité 04 : Synthèse à trous}} \hfill \badge[emerald]{Bilan}},
  coltitle=primary,colbacktitle=primarylight,titlerule=0.5pt,titlerule style=cardborder
]
{\small
\textbf{Consigne :} Complète le texte avec les mots : \textbf{malades}, \textbf{explique}, \textbf{se débarrasser}, \textbf{santé}.

\vspace{3pt}
Dans ce texte, l'auteur \drule{2.8cm} l'importance de se laver les mains. On se lave les mains pour \drule{3.2cm} des saletés afin d'éviter d'être \drule{2.8cm} et de préserver notre \drule{2.8cm}.
}
\end{tcolorbox}

\end{document}
